\documentclass[11pt,a4paper]{article}
\usepackage[T1]{fontenc}
\usepackage[utf8]{inputenc}
\usepackage{lmodern}
\usepackage{amsmath,amssymb,amsthm}
\usepackage{booktabs}
\usepackage{xcolor}
\usepackage[margin=27mm]{geometry}
\usepackage[colorlinks=true,urlcolor=blue,citecolor=blue,linkcolor=blue]{hyperref}
\newtheorem{theorem}{Theorem}
\newtheorem{proposition}[theorem]{Proposition}
\newtheorem{lemma}[theorem]{Lemma}
\newtheorem{corollary}[theorem]{Corollary}
\newtheorem{conjecture}[theorem]{Conjecture}
\theoremstyle{definition}
\newtheorem{definition}[theorem]{Definition}
\theoremstyle{plain}
\DeclareMathOperator{\Tr}{Tr}
\DeclareMathOperator{\artanh}{artanh}
\newcommand{\ket}[1]{\lvert#1\rangle}
\newcommand{\bra}[1]{\langle#1\rvert}
\newcommand{\braket}[2]{\langle#1\vert#2\rangle}
\newcommand{\norm}[1]{\lVert#1\rVert}
\newcommand{\one}{\mathbf{1}}
\newcommand{\R}{\mathbb{R}}
\newcommand{\C}{\mathbb{C}}
\newcommand{\T}{\mathsf{T}}
\newcommand{\CF}{C_{\sqrt F}}
\setlength{\emergencystretch}{2em}
\setlength{\parindent}{0pt}
\setlength{\parskip}{0.3em plus 0.1em minus 0.1em}
\title{The root-fidelity matrix bounds the Holevo quantity\\
of any three qubit states}
\author{Fable \& Astra\\[0.35em]
\normalsize Research led by Jean-Philippe Bruneton\\[-0.1em]
\small Universit\'e Paris Cit\'e, France}
\date{19 November 2026}

\begin{document}
\maketitle
\begin{abstract}
A sender who encodes a classical label into one of several quantum states can
transmit at most the Holevo quantity $\chi$ of the ensemble, and this quantity
is small when the states are hard to tell apart. Fannes, de Melo, Roga and
\.Zyczkowski proposed to bound $\chi$ for three states by a quantity built only
from the prior probabilities $p_i$ and the pairwise fidelities $F_{ij}$: the
von Neumann entropy of the $3\times3$ matrix with entries
$\sqrt{p_ip_j}\sqrt{F_{ij}}$. They proved the conjecture for pure states, and
it has since been established for special families of mixed states. We prove it
for any three qubit states, pure or mixed, with arbitrary probabilities. The
proof goes through a stronger statement valid for any number $K$ of qubit
states: $\chi$ is at most the entropy of a qubit whose determinant equals
$e_2=\sum_{i<j}p_ip_j(1-F_{ij})$, with equality exactly when all states are
pure or all are equal. We also give a sufficient condition in every dimension:
the conjecture holds whenever the three states are the image, under a positive
trace-preserving map, of three pure states with the same fidelities. This
condition covers pure states, the commuting-type families known so far, and
qubits, where the relevant map is positive but not completely positive. The
conjecture remains open in dimension three and above; a numerical search for
counterexamples in dimensions $3$, $4$ and $5$ found none.
\end{abstract}

\section{Introduction}
A sender draws a label $i$ with probability $p_i$ and prepares the quantum
state $\rho_i$. However the receiver measures, the mutual information between
the label and the outcome is bounded by the Holevo quantity
\begin{equation}\label{eq:chi}
\chi(\mathcal E)=S\Bigl(\sum_ip_i\rho_i\Bigr)-\sum_ip_iS(\rho_i),
\qquad S(\sigma)=-\Tr\sigma\log\sigma ,
\end{equation}
of the ensemble $\mathcal E=\{(p_i,\rho_i)\}$. The Holevo quantity measures how
distinguishable the states are as a family: identical states carry no
information, orthogonal pure states carry the Shannon entropy of the labels.
Fannes, de Melo, Roga and \.Zyczkowski \cite{FdMRZ2012} asked how far $\chi$ can
be controlled by pairwise data alone, namely by the probabilities and the
Uhlmann fidelities
\begin{equation}\label{eq:fid}
F_{ij}=\Bigl(\Tr\bigl|\sqrt{\rho_i}\sqrt{\rho_j}\bigr|\Bigr)^2
\end{equation}
\cite{Uhlmann1976,Jozsa1994}. For two states the answer is a theorem of Roga,
Fannes and \.Zyczkowski \cite{RFZ2010}: $\chi$ is at most the entropy of the
$2\times2$ matrix with entries $\sqrt{p_ip_j}\sqrt{F_{ij}}$. For three states,
the natural analogue is the real symmetric matrix
\begin{equation}\label{eq:C}
\CF(\mathcal E_3)=\Bigl[\sqrt{p_ip_j}\,\sqrt{F_{ij}}\Bigr]_{i,j=1}^3 ,
\end{equation}
which has diagonal $(p_1,p_2,p_3)$ and trace one. The authors of
\cite{FdMRZ2012} show that this matrix is positive semidefinite for every
triple of states, so that its entropy is defined (positivity fails in general
for four states), and they conjecture the following bound.

\begin{conjecture}[{\cite[Conjecture~1]{FdMRZ2012}}]\label{conj}
For every ensemble of three states $\{(p_i,\rho_i)\}_{i=1}^3$ on $\C^d$,
$\chi(\mathcal E_3)\le S\bigl(\CF(\mathcal E_3)\bigr)$.
\end{conjecture}

The difficulty is that $\CF$ is in general the Gram matrix of no family of
purifications of the $\rho_i$. Each fidelity is attained by a pair of
purifications, but the three optimal pairs need not be compatible, and
\cite{FdMRZ2012} observe numerically that $S(\CF)$ can lie strictly below the
smallest entropy of a correlation matrix obtained from joint purifications.
The conjecture is therefore stronger than the bound obtained from strong
subadditivity in \cite{RFZ2010}.

This note proves the conjecture for any three qubit states. The simplest
non-trivial example already shows what is measured: for three pure qubit
states, $\chi$ equals the entropy of the average state, whose determinant is
$e_2=\sum_{i<j}p_ip_j(1-F_{ij})$. Theorem~\ref{thm:qubit} shows that this
determinant controls $\chi$ for mixed qubit states as well: for any number $K$
of qubit states,
\[
\chi\le h\Bigl(\tfrac{1+\sqrt{1-4e_2}}{2}\Bigr),\qquad
h(x)=-x\log x-(1-x)\log(1-x),
\]
with equality exactly when all states are pure or all are equal. For $K=3$,
$e_2$ is the second elementary symmetric function of the eigenvalues of $\CF$,
and an integral formula for the entropy of a $3\times3$ state
(Lemma~\ref{lem:int}) shows that the right-hand side is at most $S(\CF)$
(Corollary~\ref{cor:conj}). In arbitrary dimension,
Proposition~\ref{prop:crit} reduces the conjecture to the existence of a
positive trace-preserving map sending a pure triple with the same fidelities
onto the given triple. For qubits this map exists but is not completely
positive, which explains why arguments based on purifications and strong
subadditivity fall short.

\section{Related work}\label{sec:related}
The two-state bound is due to Roga, Fannes and \.Zyczkowski \cite{RFZ2010}, who
also derived bounds on $\chi$ from the entropy of Gram matrices of joint
purifications via strong subadditivity. Conjecture~\ref{conj} is stated in
\cite{FdMRZ2012}, together with three results: the positivity of $\CF$ for three
states (obtained with D.~Vanpeteghem), the proof for three pure states, and the
weaker bound $\chi\le S(\CF\circ T_b)$ in which the off-diagonal entries are
multiplied by $b\le1/2$ in general and $b\le1/\sqrt3$ for qubits. Their
numerical support consists of $10\,000$ random ensembles drawn from the
Hilbert--Schmidt measure for $d=2,3,5,7$. Roga's thesis \cite{Roga2011}
states that the conjecture had been proved for pure states only, and adds the
case of three qubit states of which two are pure and one is mixed, with equal
probabilities $\tfrac13$. Corollary~\ref{cor:conj} contains this case and
extends it to arbitrary probabilities and arbitrary mixing.

Zhang, Wu and Fei \cite{ZWF2012} proved the conjecture under a cocycle
condition: if $U$, $V$, $W$ are the unitary parts of the polar decompositions
of $\sqrt{\rho_1}\sqrt{\rho_2}$, $\sqrt{\rho_1}\sqrt{\rho_3}$ and
$\sqrt{\rho_2}\sqrt{\rho_3}$, and $V=UW$, then $\chi\le S(\CF)$. The condition
says that the three pairwise optimal purifications are mutually compatible; it
holds in particular for commuting states. Proposition~\ref{prop:crit} recovers
this case with the partial trace as the positive map.

Among the later works citing \cite{FdMRZ2012}, Shirokov \cite{Shirokov2025}
derives upper bounds on the Holevo quantity from an entropic inequality
without returning to Conjecture~\ref{conj}, and Vuong and Osaka
\cite{VuongOsaka2026} cite \cite{FdMRZ2012} for the positivity of the
root-fidelity matrix in a study of three-point trace inequalities, without
discussing the Holevo quantity. In the documented search described in
Section~\ref{sec:scope}, no proof of the general qubit case and no
counterexample were found.

The single external tool used below is the monotonicity of the relative
entropy under positive, not necessarily completely positive, trace-preserving
maps, proved by M\"uller-Hermes and Reeb \cite{MHR2017}. The fidelity formula
for qubits and the integral representation of the logarithm are classical.

\section{Two lemmas on entropies}
Entropies use the natural logarithm; none of the statements depends on the
base. For a positive semidefinite $3\times3$ matrix $\sigma$ of unit trace
with eigenvalues $\lambda_1,\lambda_2,\lambda_3$, write
$e_2(\sigma)=\lambda_1\lambda_2+\lambda_1\lambda_3+\lambda_2\lambda_3$, the sum
of the principal $2\times2$ minors, and $e_3(\sigma)=\det\sigma$. For
$\sigma=\CF$ one has
\begin{equation}\label{eq:e2e3}
e_2=\sum_{i<j}p_ip_j(1-F_{ij}),\qquad
e_3=p_1p_2p_3\Bigl(1-F_{12}-F_{13}-F_{23}+2\sqrt{F_{12}F_{13}F_{23}}\Bigr).
\end{equation}

\begin{lemma}\label{lem:int}
For every such $\sigma$,
\[
S(\sigma)=\int_0^\infty\Bigl(\frac{t^2+2e_2t+3e_3}{t^3+t^2+e_2t+e_3}-\frac1{1+t}\Bigr)dt .
\]
Consequently, if $\sigma$ and $\sigma'$ have the same $e_2$ and
$e_3(\sigma)\le e_3(\sigma')$, then $S(\sigma)\le S(\sigma')$.
\end{lemma}

\begin{proof}
For $\lambda>0$, $\log\lambda=\int_0^\infty\bigl(\frac1{1+t}-\frac1{\lambda+t}\bigr)dt$,
so that $-\lambda\log\lambda=\int_0^\infty\bigl(\frac{\lambda}{\lambda+t}-\frac{\lambda}{1+t}\bigr)dt$;
the integrand equals $\lambda(1-\lambda)/((\lambda+t)(1+t))$ and is integrable,
and the identity also holds for $\lambda=0$. Summing over the three eigenvalues,
whose sum is one,
$S(\sigma)=\int_0^\infty\bigl(\sum_k\frac{\lambda_k}{\lambda_k+t}-\frac1{1+t}\bigr)dt$.
With $P(t)=\prod_k(t+\lambda_k)=t^3+t^2+e_2t+e_3$ one has
$\sum_k\frac{\lambda_k}{\lambda_k+t}=3-tP'(t)/P(t)=(t^2+2e_2t+3e_3)/P(t)$, which
is the formula. For fixed $t>0$ and $e_2$, the derivative of
$x\mapsto(t^2+2e_2t+3x)/(t^3+t^2+e_2t+x)$ is
$(3t^3+2t^2+e_2t)/(t^3+t^2+e_2t+x)^2>0$, so the integrand increases pointwise
with $e_3$.
\end{proof}

Lemma~\ref{lem:int} recovers the pure case of \cite{FdMRZ2012}. For pure
states $\rho_i=\ket{\varphi_i}\bra{\varphi_i}$, the Gram matrix
$G=[\sqrt{p_ip_j}\braket{\varphi_i}{\varphi_j}]$ has the same nonzero spectrum
as $\sum_ip_i\rho_i$, so $\chi=S(G)$, and $\CF=|G|$ entrywise. The two
matrices have the same $e_2$, and
$\det|G|-\det G=2\bigl(|G_{12}G_{23}G_{31}|-\operatorname{Re}G_{12}G_{23}G_{31}\bigr)\ge0$,
hence $\chi=S(G)\le S(\CF)$.

\begin{lemma}\label{lem:qubit}
Let $f(D)=h\bigl(\tfrac{1+\sqrt{1-4D}}{2}\bigr)$ be the entropy of a qubit
state of determinant $D\in[0,\tfrac14]$, and $\phi(s)=f(s^2/4)$ for
$s\in[0,1]$. Then $f$ is increasing, strictly concave and $f(0)=0$; hence it
is subadditive, $f(x+y)\le f(x)+f(y)$ for $x,y\ge0$ with $x+y\le\tfrac14$, with
equality only if $x=0$ or $y=0$. The function $\phi$ is strictly convex.
\end{lemma}

\begin{proof}
Let $g(r)=h\bigl(\tfrac{1+r}{2}\bigr)$, so $g'(r)=-\artanh r$. With
$r=\sqrt{1-4D}$ one finds $f'(D)=2\artanh(r)/r=2\sum_{k\ge0}r^{2k}/(2k+1)$,
which is positive and strictly increasing in $r$, hence strictly decreasing in
$D$; so $f$ is increasing and strictly concave. Concavity and $f(0)=0$ give
$f(x)\ge\frac{x}{x+y}f(x+y)$ and $f(y)\ge\frac{y}{x+y}f(x+y)$, with equality in
either only when the corresponding variable vanishes; summing gives
subadditivity and its equality case. With $r=\sqrt{1-s^2}$ one has
$\phi'(s)=\psi(r):=\sqrt{1-r^2}\,\artanh(r)/r$ and, for $0<r<1$,
\[
\frac{d}{dr}\log\psi(r)=-\frac{1}{r(1-r^2)}+\frac{1}{(1-r^2)\artanh r}<0,
\]
because $\artanh r>r$. Thus $\psi$ is strictly decreasing in $r$, while $r$ is
strictly decreasing in $s$, so $\phi'$ is strictly increasing and $\phi$ is
strictly convex.
\end{proof}

\section{Qubit ensembles}
\begin{theorem}\label{thm:qubit}
Let $\{(p_i,\rho_i)\}_{i=1}^K$ be an ensemble of $K\ge2$ qubit states, pure or
mixed, with arbitrary probabilities, and let
$e_2=\sum_{i<j}p_ip_j(1-F_{ij})$. Then $e_2\le\tfrac14$ and
\[
\chi\le f(e_2)=h\Bigl(\frac{1+\sqrt{1-4e_2}}{2}\Bigr).
\]
Equality holds if and only if all the states are pure or all the states are
equal.
\end{theorem}

\begin{proof}
Write $\rho_i=\tfrac12(\one+\vec a_i\cdot\vec\sigma)$ with Bloch vector
$\vec a_i$ of length $r_i\le1$, set $s_i=\sqrt{1-r_i^2}=2\sqrt{\det\rho_i}$, and
attach to $\rho_i$ the unit vector $\hat a_i=(\vec a_i,s_i)\in\R^4$. For a
positive $2\times2$ matrix $M$, $(\Tr\sqrt M)^2=\Tr M+2\sqrt{\det M}$; applied
to $M=\sqrt{\rho_i}\rho_j\sqrt{\rho_i}$, whose square-root trace is
$\Tr|\sqrt{\rho_i}\sqrt{\rho_j}|$, this gives
$F_{ij}=\Tr\rho_i\rho_j+2\sqrt{\det\rho_i\det\rho_j}
=\tfrac12(1+\vec a_i\cdot\vec a_j+s_is_j)$, that is,
\begin{equation}\label{eq:qubitfid}
F_{ij}=\frac{1+\langle\hat a_i,\hat a_j\rangle}{2}.
\end{equation}
Let $m=\bigl|\sum_ip_i\vec a_i\bigr|$, $\bar s=\sum_ip_is_i$ and
$R=\bigl\|\sum_ip_i\hat a_i\bigr\|=\sqrt{m^2+\bar s^2}\le1$. Since
$e_2=\tfrac12\sum_{i,j}p_ip_j(1-F_{ij})$, equation~\eqref{eq:qubitfid} gives
$R^2=\sum_{i,j}p_ip_j(2F_{ij}-1)=1-4e_2$, so
$e_2=\tfrac14(1-m^2-\bar s^2)\in[0,\tfrac14]$.

The entropy of $\rho_i$ is $\phi(s_i)$ and that of the average state, of Bloch
vector length $m$, is $f\bigl(\tfrac{1-m^2}{4}\bigr)$. By convexity of $\phi$
(Lemma~\ref{lem:qubit}), $\sum_ip_i\phi(s_i)\ge\phi(\bar s)=f(\bar s^2/4)$. By
subadditivity of $f$ applied to $x=e_2$ and $y=\bar s^2/4$, whose sum is
$(1-m^2)/4$,
$f\bigl(\tfrac{1-m^2}{4}\bigr)\le f(e_2)+f(\bar s^2/4)$. Therefore
\[
\chi=f\Bigl(\frac{1-m^2}{4}\Bigr)-\sum_ip_i\phi(s_i)
\le f(e_2)+f\Bigl(\frac{\bar s^2}{4}\Bigr)-f\Bigl(\frac{\bar s^2}{4}\Bigr)=f(e_2).
\]
Equality forces equality in the subadditivity step, hence $\bar s=0$ or
$e_2=0$. If $\bar s=0$, all states are pure, and then
$\chi=f\bigl(\tfrac{1-m^2}{4}\bigr)=f(e_2)$. If $e_2=0$, then $R=1$, all
$\hat a_i$ coincide, the states are equal and $\chi=0=f(0)$. Conversely, in
both cases equality holds.
\end{proof}

\begin{corollary}\label{cor:conj}
Conjecture~\ref{conj} holds for every ensemble of three qubit states:
$\chi\le f(e_2)\le S(\CF)$. The second inequality is strict whenever
$\det\CF>0$.
\end{corollary}

\begin{proof}
The matrix $\CF$ is positive semidefinite \cite{FdMRZ2012} with unit trace,
second invariant $e_2$ and determinant $e_3\ge0$. The spectrum
$(\lambda,1-\lambda,0)$ with $\lambda(1-\lambda)=e_2$, which exists since
$e_2\le\tfrac14$, has invariants $(e_2,0)$ and entropy $f(e_2)$.
Lemma~\ref{lem:int} gives $f(e_2)\le S(\CF)$, strictly if $e_3>0$ because the
integrand is then strictly larger for every $t>0$, and
Theorem~\ref{thm:qubit} gives $\chi\le f(e_2)$.
\end{proof}

Theorem~\ref{thm:qubit} isolates the mechanism. For qubits, $\chi$ is
controlled by the single invariant $e_2$, equivalently by the average fidelity
$\sum_{i,j}p_ip_jF_{ij}=1-2e_2$, and the determinant of $\CF$ only loosens the
bound. For $K=2$, $e_2=p_1p_2(1-F_{12})=\det\CF$ and Theorem~\ref{thm:qubit}
reduces to the qubit case of the two-state bound of \cite{RFZ2010}. For
$K\ge4$ the bound $\chi\le f(e_2)$ remains valid although the root-fidelity
matrix may fail to be positive.

\section{A sufficient criterion in any dimension}
\begin{proposition}\label{prop:crit}
Let $\{(p_i,\rho_i)\}_{i=1}^3$ be an ensemble on $\C^d$. Suppose there exist
unit vectors $\phi_1,\phi_2,\phi_3\in\C^n$ and a positive trace-preserving
linear map $\Lambda:M_n(\C)\to M_d(\C)$ such that
$\Lambda(\ket{\phi_i}\bra{\phi_i})=\rho_i$ and
$|\braket{\phi_i}{\phi_j}|^2=F_{ij}$ for all $i,j$. Then
$\chi\le S(\CF)$.
\end{proposition}

\begin{proof}
Write $\bar\rho=\sum_ip_i\rho_i$ and
$\bar\phi=\sum_ip_i\ket{\phi_i}\bra{\phi_i}$, so that $\bar\rho=\Lambda(\bar\phi)$.
With the relative entropy $D(\rho\|\sigma)=\Tr\rho(\log\rho-\log\sigma)$, one
has $\chi=\sum_ip_iD(\rho_i\|\bar\rho)$. Monotonicity of the relative entropy
under positive trace-preserving maps \cite{MHR2017} gives
$\chi\le\sum_ip_iD(\phi_i\|\bar\phi)=S(\bar\phi)=S(G_\phi)$, where
$G_\phi=[\sqrt{p_ip_j}\braket{\phi_i}{\phi_j}]$. Since $|G_\phi|=\CF$
entrywise, the pure case following Lemma~\ref{lem:int} gives
$S(G_\phi)\le S(\CF)$.
\end{proof}

The three families for which the conjecture was known satisfy this criterion.
For pure states, $\Lambda$ is the identity. Under the cocycle condition of
Zhang, Wu and Fei \cite{ZWF2012}, in particular for commuting states, the
$\phi_i$ can be taken to be mutually compatible optimal purifications and
$\Lambda$ the partial trace, which is completely positive.

For three qubit states the criterion holds with a map that is positive but
generally not completely positive. The vectors $\hat a_i\in\R^4$ of the proof of
Theorem~\ref{thm:qubit} span a subspace of dimension at most three, so there
are unit vectors $\hat b_i\in\R^3$ with the same Gram matrix. They are Bloch
vectors of pure qubit states $\phi_i$ with
$|\braket{\phi_i}{\phi_j}|^2=\tfrac12(1+\langle\hat b_i,\hat b_j\rangle)=F_{ij}$
by~\eqref{eq:qubitfid}. Let $O$ be a partial isometry sending each $\hat b_i$
to $\hat a_i$, which exists because the Gram matrices agree, let
$\Pi:\R^4\to\R^3$ forget the fourth coordinate, and let $\Lambda$ be the unital
qubit map with Bloch matrix $M=\Pi O$, that is,
$\Lambda\bigl(\tfrac12(\one+\vec x\cdot\vec\sigma)\bigr)=\tfrac12(\one+M\vec x\cdot\vec\sigma)$.
Then $M\hat b_i=\vec a_i$ and $\norm{M}\le1$, so $\Lambda$ maps the Bloch ball into
itself: it is positive and trace preserving. For pure qubit states $S(G_\phi)$
is the entropy of the average state, whose Bloch vector has squared length
$\|\sum_ip_i\hat b_i\|^2=R^2=1-4e_2$, so $S(G_\phi)=f(e_2)$ and
Proposition~\ref{prop:crit} gives back the chain of Corollary~\ref{cor:conj}.

When the $\hat a_i$ are linearly independent and at least one state is mixed,
$\Lambda$ is not completely positive. In that case $O$ is an isometry from
$\R^3$ into $\R^4$, and $M^\T M=O^\T(\one-e_4e_4^\T)O=\one-vv^\T$ with
$v=O^\T e_4$; since $\langle v,\hat b_i\rangle=s_i$, $v\ne0$. The singular
values of $M$ are therefore $(1,1,\sigma_3)$ with
$\sigma_3=\sqrt{1-\norm v^2}<1$. Write $M=O_1DO_2$ with $O_1,O_2\in SO(3)$ and
$D=\mathrm{diag}(\lambda_1,\lambda_2,\lambda_3)$, $|\lambda_k|$ the singular
values. Rotations of the Bloch ball are unitary conjugations, so $\Lambda$ is
completely positive if and only if the map with Bloch matrix $D$ is. That map
is $X\mapsto\sum_{k=0}^3q_k\sigma_kX\sigma_k$ with $\sigma_0=\one$ and
$4q_0=1+\lambda_1+\lambda_2+\lambda_3$,
$4q_1=1+\lambda_1-\lambda_2-\lambda_3$,
$4q_2=1-\lambda_1+\lambda_2-\lambda_3$,
$4q_3=1-\lambda_1-\lambda_2+\lambda_3$; its Choi matrix has eigenvalues
$2q_k$, so it is completely positive exactly when all $q_k\ge0$. The sign
patterns of $(\lambda_1,\lambda_2)$ in these four expressions run over all four
choices, so one of them equals $1-|\lambda_1|-|\lambda_2|\pm\lambda_3=-1\pm\lambda_3<0$.
This explains why bounds derived from joint purifications and strong
subadditivity, which only use completely positive maps, do not reach the
conjecture in the qubit case.

\section{Dimension three and higher}
For $d\ge3$ the conjecture remains open. The proof of
Theorem~\ref{thm:qubit} rests on two features specific to qubits: the
fidelity~\eqref{eq:qubitfid} is an affine function of an inner product, and the
entropy of a qubit state depends only on its determinant. Neither survives in
dimension three, where $\chi$ can exceed $\log2$ and the bound by $e_2$ alone
fails: three orthogonal pure states with equal probabilities give
$e_2=\tfrac13>\tfrac14$ and $\chi=\log3$. Proposition~\ref{prop:crit} suggests a
route: show that every triple of states is the image, under a positive
trace-preserving map, of a pure triple with the same fidelities, the phases of
the pure triple being free. This is established above for pure states, for
triples satisfying the cocycle condition and for qubits; it has been neither
proved nor refuted in general, and no numerical test of this existence
statement is reported here.

During the preparation of this note we searched numerically for a
counterexample, maximising $\Delta=\chi-S(\CF)$ by gradient ascent in double
precision. With free states of prescribed rank and free probabilities in
dimensions $3$ and $4$, maximising $\chi/S(\CF)$, the runs converge to
configurations of equality $\Delta=0$, several with $\chi$ between $0.32$ and
$0.66$; the best double-precision values of $\Delta$, of order $10^{-14}$,
become negative when re-evaluated with $40$ significant digits. On the unitary
orbits of fixed spectra with fixed probabilities, the largest ratio
$\chi/S(\CF)$ found was $0.99991$ for $d=2$, $0.99977$ for $d=3$, $0.99384$ for
$d=4$ and $0.99850$ for $d=5$, with largest $\Delta=-2.8\times10^{-5}$.
Full-rank perturbations of low-rank equality points, of size $10^{-1}$ to
$10^{-3}$ in dimensions $3$, $4$ and $5$, made $\Delta$ strictly negative, and
reoptimisation brought it back towards $0$ from below. These runs used a
separate script and are not part of the public bundle; the public checker
contains a reduced, fixed-seed version (Section~\ref{sec:verification}). The
search is local and does not constitute evidence beyond the configurations it
visits.

\section{Verification}\label{sec:verification}
The results above are proved analytically. The public
\href{https://github.com/jpbruneton/open_quantum_problems/blob/main/research/entries/fidelity-matrix-holevo-bound/checks/check_fidelity_holevo.py}{Python checker}
tests their ingredients in floating-point arithmetic and writes its
\href{https://github.com/jpbruneton/open_quantum_problems/blob/main/research/entries/fidelity-matrix-holevo-bound/data/verification.json}{verification record}
to the same entry. It checks the integral formula of Lemma~\ref{lem:int} on
$20$ spectra (largest error below $10^{-15}$) and the inequality
$S(|G|)\ge S(G)$ on $5000$ pure triples; the concavity of $f$ and convexity of
$\phi$ on grids of $20\,001$ points; formula~\eqref{eq:qubitfid} and
$e_2=(1-R^2)/4$ against matrix computations of the fidelity; the two
inequalities $\chi\le f(e_2)\le S(\CF)$ on $10\,000$ random three-qubit
ensembles, including nearly pure, exactly pure and nearly maximally mixed
states, where the smallest slack is of order $-10^{-15}$, at rounding level on
equality cases; the inequality $\chi\le f(e_2)$ for $K=2,4,5,6$ on $20\,000$
ensembles and under $20$ Nelder--Mead minimisations of the slack; and the
construction of Section~5 on $500$ triples, where $\norm M\le1$, the fidelities
are preserved and $\Lambda$ fails to be completely positive in every case. For
$d=3,4,5$ it evaluates $\Delta$ on $4500$ random ensembles and runs $24$ short
fixed-seed maximisations of $\Delta$, over free ensembles and over unitary
orbits; all values found are negative. The checker requires NumPy and SciPy
and runs in under a minute.

\section{Scope and limitations}\label{sec:scope}
\emph{Analytic proof.} Lemmas~\ref{lem:int} and~\ref{lem:qubit},
Theorem~\ref{thm:qubit} with its equality case, Corollary~\ref{cor:conj}
(Conjecture~\ref{conj} for any three qubit states) and
Proposition~\ref{prop:crit}, the last relying on the published theorem of
M\"uller-Hermes and Reeb \cite{MHR2017}, and on the positivity of $\CF$ proved
in \cite{FdMRZ2012}.

\emph{Numerical evidence only.} The absence of counterexamples in dimensions
$3$, $4$ and $5$. The search is adversarial but local and does not explore
dimensions above $5$.

\emph{Open.} Conjecture~\ref{conj} for $d\ge3$, and the existence, for every
triple of states, of a pure triple with the same fidelities mapped onto it by a
positive trace-preserving map.

\emph{Bibliographic priority.} The source \cite{FdMRZ2012} and Roga's thesis
\cite{Roga2011} were read in their TeX sources and \cite{ZWF2012} in full. The
works citing \cite{FdMRZ2012} were listed through the Semantic Scholar
interface ($24$ citations between 2011 and 2026) and complemented by three web
searches; \cite{Shirokov2025} and \cite{VuongOsaka2026} were read in full, the
other citing works only by title. Theorem~\ref{thm:qubit} and
Corollary~\ref{cor:conj} were not found in this documented search, which does
not establish priority. The qubit fidelity formula, the integral representation
of the logarithm and the convexity of $\phi$, which is the same function that
expresses the entanglement of formation in terms of the concurrence, are
classical.

\clearpage
\begingroup\small
\begin{thebibliography}{99}
\setlength{\itemsep}{2pt}
\bibitem{FdMRZ2012}
M. Fannes, F. de Melo, W. Roga and K. \.Zyczkowski,
\emph{Matrices of fidelities for ensembles of quantum states and the Holevo
quantity},
Quantum Inf. Comput. \textbf{12}, 472--489 (2012);
\href{https://doi.org/10.26421/QIC12.5-6-8}{doi:10.26421/QIC12.5-6-8};
\href{https://arxiv.org/abs/1104.2271}{arXiv:1104.2271}.

\bibitem{RFZ2010}
W. Roga, M. Fannes and K. \.Zyczkowski,
\emph{Universal Bounds for the Holevo Quantity, Coherent Information, and the
Jensen-Shannon Divergence},
Phys. Rev. Lett. \textbf{105}, 040505 (2010);
\href{https://doi.org/10.1103/PhysRevLett.105.040505}{doi:10.1103/PhysRevLett.105.040505}.

\bibitem{Roga2011}
W. Roga,
\emph{Entropy of quantum channel in the theory of quantum information},
PhD thesis (2011);
\href{https://arxiv.org/abs/1108.5065}{arXiv:1108.5065}.

\bibitem{ZWF2012}
L. Zhang, J. Wu and S.-M. Fei,
\emph{Universal upper bound for the Holevo information induced by a quantum
operation},
Phys. Lett. A \textbf{376}, 3588--3592 (2012);
\href{https://doi.org/10.1016/j.physleta.2012.10.014}{doi:10.1016/j.physleta.2012.10.014};
\href{https://arxiv.org/abs/1110.5979}{arXiv:1110.5979}.

\bibitem{MHR2017}
A. M\"uller-Hermes and D. Reeb,
\emph{Monotonicity of the Quantum Relative Entropy Under Positive Maps},
Ann. Henri Poincar\'e \textbf{18}, 1777--1788 (2017);
\href{https://doi.org/10.1007/s00023-017-0550-9}{doi:10.1007/s00023-017-0550-9};
\href{https://arxiv.org/abs/1512.06117}{arXiv:1512.06117}.

\bibitem{Uhlmann1976}
A. Uhlmann,
\emph{The ``transition probability'' in the state space of a $*$-algebra},
Rep. Math. Phys. \textbf{9}, 273--279 (1976);
\href{https://doi.org/10.1016/0034-4877(76)90060-4}{doi:10.1016/0034-4877(76)90060-4}.

\bibitem{Jozsa1994}
R. Jozsa,
\emph{Fidelity for Mixed Quantum States},
J. Mod. Opt. \textbf{41}, 2315--2323 (1994);
\href{https://doi.org/10.1080/09500349414552171}{doi:10.1080/09500349414552171}.

\bibitem{Shirokov2025}
M. E. Shirokov,
\emph{Upper bounds on the Holevo quantity arising from the fundamental entropic
inequality}, preprint (2025);
\href{https://arxiv.org/abs/2506.05335}{arXiv:2506.05335}.

\bibitem{VuongOsaka2026}
T. D. Vuong and H. Osaka,
\emph{Three-Point Trace Positivity: Schatten Norm Compression and Interpolating
Metrics}, preprint (2026);
\href{https://arxiv.org/abs/2610.11021}{arXiv:2610.11021}.
\end{thebibliography}
\endgroup
\end{document}
